\begin{ex}
Dùng thước dây đo chiều dài bảng lớp được 3 mét. Đổi 3 mét ra đề-xi-mét là bao nhiêu đề-xi-mét?
\shortans{30}
\loigiai{
$3\text{ m} = 30\text{ dm}$.
}
\end{ex}

\begin{ex}
Đo chiều dài lớp học được 8 mét. Đổi 8 mét ra đề-xi-mét là bao nhiêu đề-xi-mét?
\shortans{80}
\loigiai{
$8\text{ m} = 80\text{ dm}$.
}
\end{ex}

\begin{ex}
Chiều dài cuốn sách Toán 2 đo được 24 cm. Chiều rộng đo được 17 cm. Chiều dài hơn chiều rộng bao nhiêu xăng-ti-mét?
\shortans{7}
\loigiai{
Hơn: $24 - 17 = 7$ cm.
}
\end{ex}

\begin{ex}
Bạn Lan bước 5 bước chân, mỗi bước chân dài khoảng 40 cm. Quãng đường Lan bước dài bao nhiêu xăng-ti-mét?
\shortans{200}
\loigiai{
Dài: $40 \times 5 = 200$ cm.
}
\end{ex}

\begin{ex}
Đổi 200 cm ra mét được bao nhiêu mét?
\shortans{2}
\loigiai{
$200\text{ cm} = 2\text{ m}$.
}
\end{ex}

\begin{ex}
Một cái bàn học dài 120 cm. Đổi ra đề-xi-mét là bao nhiêu đề-xi-mét?
\shortans{12}
\loigiai{
$120\text{ cm} = 12\text{ dm}$.
}
\end{ex}

\begin{ex}
Một thanh nhôm dài 2 m, bác thợ cắt đi 8 dm. Thanh nhôm còn lại dài bao nhiêu đề-xi-mét?
\shortans{12}
\loigiai{
$2\text{ m} = 20\text{ dm}$. Còn lại: $20 - 8 = 12$ dm.
}
\end{ex}

\begin{ex}
Cây cau cao 9 m, cây chuối cao 3 m. Cây cau cao gấp mấy lần cây chuối? Cây cau cao hơn cây chuối bao nhiêu mét?
\shortans{6}
\loigiai{
Cao hơn: $9 - 3 = 6$ mét.
}
\end{ex}

\begin{ex}
Đoạn thẳng AB dài 15 cm, vẽ thêm đoạn BC dài 25 cm thẳng hàng với AB. Cả đoạn thẳng AC dài bao nhiêu xăng-ti-mét?
\shortans{40}
\loigiai{
Đoạn AC dài: $15 + 25 = 40$ cm.
}
\end{ex}

\begin{ex}
Đoạn đường quanh sân trường hình chữ nhật có chu vi là 300 m. Bạn Hùng chạy 2 vòng quanh sân trường thì chạy được bao nhiêu mét?
\shortans{600}
\loigiai{
Chạy được: $300 \times 2 = 600$ m.
}
\end{ex}

\begin{ex}
Một cuộn chỉ dài 100 m, người ta dùng hết 65 m. Cuộn chỉ còn lại bao nhiêu mét?
\shortans{35}
\loigiai{
Còn lại: $100 - 65 = 35$ m.
}
\end{ex}

\begin{ex}
Bạn An cao 128 cm, bạn Bình cao 135 cm. Bạn Bình cao hơn bạn An bao nhiêu xăng-ti-mét?
\shortans{7}
\loigiai{
Cao hơn: $135 - 128 = 7$ cm.
}
\end{ex}

\begin{ex}
Bác nông dân rào vườn hết 85 m lưới thép, bác mua cuộn lưới dài 100 m. Bác còn thừa bao nhiêu mét lưới thép?
\shortans{15}
\loigiai{
Thừa: $100 - 85 = 15$ m.
}
\end{ex}

\begin{ex}
Một chiếc cột điện cao 12 m, chiếc cột đèn cao 8 m. Cột điện cao hơn cột đèn bao nhiêu mét?
\shortans{4}
\loigiai{
Cao hơn: $12 - 8 = 4$ m.
}
\end{ex}

\begin{ex}
Đoạn đường từ cổng trường đến thư viện dài 45 m, từ thư viện đến nhà ăn dài 35 m. Quãng đường từ cổng trường đến nhà ăn đi qua thư viện dài bao nhiêu mét?
\shortans{80}
\loigiai{
Dài: $45 + 35 = 80$ m.
}
\end{ex}

\begin{ex}
Một con sâu đo một chiếc lá dài 12 cm, mỗi lần sâu vươn mình tiến được 2 cm. Sâu cần đo bao nhiêu lần để đi hết chiếc lá?
\shortans{6}
\loigiai{
Cần đo: $12 : 2 = 6$ lần.
}
\end{ex}

\begin{ex}
Đoạn thẳng MN dài 5 dm, đoạn thẳng PQ dài 30 cm. Đoạn MN dài hơn đoạn PQ bao nhiêu xăng-ti-mét?
\shortans{20}
\loigiai{
$5\text{ dm} = 50\text{ cm}$. Hơn: $50 - 30 = 20$ cm.
}
\end{ex}

\begin{ex}
Một sợi dây thừng dài 18 m được cắt làm 2 đoạn bằng nhau. Mỗi đoạn dài bao nhiêu mét?
\shortans{9}
\loigiai{
Mỗi đoạn dài: $18 : 2 = 9$ m.
}
\end{ex}

\begin{ex}
Một sợi dây đồng dài 20 dm được cắt thành các đoạn 5 dm. Cắt được bao nhiêu đoạn?
\shortans{4}
\loigiai{
Cắt được: $20 : 5 = 4$ đoạn.
}
\end{ex}

\begin{ex}
Đo chiều cao bục giảng được 20 cm, bục có 2 bậc bằng nhau. Mỗi bậc cao bao nhiêu xăng-ti-mét?
\shortans{10}
\loigiai{
Mỗi bậc cao: $20 : 2 = 10$ cm.
}
\end{ex}

\begin{ex}
Một chiếc thang có 10 bậc, khoảng cách giữa các bậc là 2 dm. Chiều dài chiếc thang tính từ bậc 1 đến bậc 10 (gồm 9 khoảng cách) là bao nhiêu đề-xi-mét?
\shortans{18}
\loigiai{
Chiều dài: $2 \times 9 = 18$ dm.
}
\end{ex}

\begin{ex}
Tính: $45\text{ m} + 35\text{ m} - 20\text{ m}$.
\shortans{60}
\loigiai{
Ta có $45 + 35 = 80$, rồi $80 - 20 = 60$ m.
}
\end{ex}

\begin{ex}
Tính: $100\text{ m} - 40\text{ m} - 30\text{ m}$.
\shortans{30}
\loigiai{
Ta có $100 - 40 = 60$, rồi $60 - 30 = 30$ m.
}
\end{ex}

\begin{ex}
Tính: $500\text{ m} + 500\text{ m}$. Kết quả đổi ra ki-lô-mét là bao nhiêu ki-lô-mét?
\shortans{1}
\loigiai{
Ta có $500 + 500 = 1000\text{ m} = 1$ km.
}
\end{ex}

\begin{ex}
Một người đi bộ 1 giờ được 4 km. Trong 2 giờ người đó đi được bao nhiêu ki-lô-mét?
\shortans{8}
\loigiai{
Đi được: $4 \times 2 = 8$ km.
}
\end{ex}
